\section*{Chapter \ref{ch:g12:nuclear-energy} --- Nuclear Energy: Fission, Fusion, $E=mc^2$}

\begin{solution}{exo:g12:nuclear-energy:1}
(a) $\qty{1}{eV} = \qty{1.60e-19}{J}$;
$\qty{1}{MeV} = \qty{1.60e-13}{J}$.
(b) $200 \times \num{1.60e-13} = \qty{3.2e-11}{J}$.
(c) $E = \num{1.66054e-27} \times \num{9.00e16} = \qty{1.49e-10}{J}
= \num{1.49e-10}/\num{1.60e-13} \approx \qty{9.3e2}{MeV}$ --- the
\qty{931.5}{MeV} of the course, up to our rounded $c$.
\end{solution}

\begin{solution}{exo:g12:nuclear-energy:2}
$E = \num{6.0e-3} \times \num{9.00e16} = \qty{5.4e14}{J}$;
$\num{5.4e14}/\num{3.6e7} = \num{1.5e7}$ days $\approx \num{4.1e4}$
years --- forty millennia on one sugar cube.
\end{solution}

\begin{solution}{exo:g12:nuclear-energy:3}
$\Delta m = 6 \times 1.00728 + 6 \times 1.00866 - 11.99671 =
\qty{0.09893}{u}$; $E_b = 0.09893 \times 931.5 \approx \qty{92.2}{MeV}$;
$E_b/A \approx \qty{7.68}{MeV}$ per \omterm{def:g11:fundamental-interactions:ladder}{nucleon}.
\end{solution}

\begin{solution}{exo:g12:nuclear-energy:4}
(a) $A = 236 - 144 - 3 = 89$, $Z = 92 - 56 = 36$:
${}^{89}_{36}\mathrm{Kr}$, krypton.
(b) $A = 236 - 97 - 2 = 137$, $Z = 92 - 40 = 52$:
${}^{137}_{52}\mathrm{Te}$, tellurium.
\end{solution}

\begin{solution}{exo:g12:nuclear-energy:5}
About $7.1$, $7.7$, $8.8$ and \qty{7.6}{MeV} per \omterm{def:g11:fundamental-interactions:ladder}{nucleon}; iron-56 is
the most bound. Uranium can release \omterm{def:g10:energy-conservation:energy}{energy} by \omterm{def:g12:nuclear-energy:fission}{fission}, helium-4 and
carbon-12 by \omterm{def:g12:nuclear-energy:fusion}{fusion} --- both moves climb toward iron; iron itself has
nowhere to go.
\end{solution}

\begin{solution}{exo:g12:nuclear-energy:6}
$\Delta m = 1.00728 + 1.00866 - 2.01355 = \qty{0.00239}{u}$, so
$E_b \approx \qty{2.23}{MeV}$ --- about a million times the few
\unit{eV} of a chemical bond.
\end{solution}

\begin{solution}{exo:g12:nuclear-energy:7}
Before: $\qty{236.00216}{u}$; after: $143.89220 + 88.89790 +
3 \times 1.00866 = \qty{235.81608}{u}$; $\Delta m = \qty{0.18608}{u}$,
so $E \approx \qty{173}{MeV} = 173 \times \num{1.60e-13} \approx
\qty{2.8e-11}{J}$.
\end{solution}

\begin{solution}{exo:g12:nuclear-energy:8}
$\Delta m = 2 \times 2.01355 - 3.01493 - 1.00866 = \qty{0.00351}{u}$:
$E \approx \qty{3.27}{MeV}$. D--T makes helium-4, anomalously bound
(\qty{7.1}{MeV} per \omterm{def:g11:fundamental-interactions:ladder}{nucleon}); helium-3 sits far lower on the curve, so
the step releases much less.
\end{solution}

\begin{solution}{exo:g12:nuclear-energy:9}
(a) The neutron is neutral: no \omterm{def:g12:nuclear-energy:barrier}{Coulomb barrier}, \omterm{def:g12:nuclear-energy:fission}{fission} fires cold.
Fusing nuclei are both positive and must climb the barrier, so only
$\sim \qty{1e8}{K}$ agitation gives collisions violent enough. (b) The
core holds an astronomical number of protons; the rare fastest ones in
the crowd do the fusing, and the Sun has billions of years to spare
--- a slow burn is exactly what a star wants. (c) There is no neutron
avalanche: each \omterm{def:g12:nuclear-energy:fusion}{fusion} is bought by its own collision, and any loss of
confinement cools the \omterm{def:g12:nuclear-energy:barrier}{plasma} and puts the fire out.
\end{solution}

\begin{solution}{exo:g12:nuclear-energy:10}
$\Delta m = 4 \times 1.00728 - 4.00151 = \qty{0.02761}{u}$, so
$E \approx \qty{25.7}{MeV}$; fraction
$0.02761/4.02912 \approx 0.7\%$ of the mass.
\end{solution}

\begin{solution}{exo:g12:nuclear-energy:11}
(a) It slows \omterm{def:g12:nuclear-energy:fission}{fission} neutrons by collisions; slow neutrons are far
more likely to \omterm{def:g12:nuclear-energy:fission}{fission} uranium-235, so the chain needs them. (b) They
absorb neutrons, holding the chain at one neutron per \omterm{def:g12:nuclear-energy:fission}{fission} ---
pushed in, it dies. (c) No water, no moderation: the neutrons stay
fast, miss, and the chain chokes --- the design fails toward
``off''.
\end{solution}

\begin{solution}{exo:g12:nuclear-energy:12}
\omterm{def:g12:nuclear-energy:fission}{Fissions} double each generation, so after $n$ generations about $2^n$
have split; $2^{n} = \num{2.56e21} \approx 2 \times (10^{3})^{7}
\approx 2^{71}$ gives $n \approx 71$. Time: $71 \times \qty{10}{ns}
\approx \qty{0.7}{\micro s}$ --- a gram of fuel, three tonnes of coal's
worth, in under a microsecond: criticality is a cliff, not a slope.
\end{solution}

\begin{solution}{exo:g12:nuclear-energy:13}
$\Delta m = 55.92066 - 2 \times 27.96924 = \qty{-0.01782}{u}$:
$E \approx \qty{-16.6}{MeV}$ --- the split \emph{absorbs} \omterm{def:g10:energy-conservation:energy}{energy}.
Iron sits at the peak: fissioning or fusing it both go downhill in
binding, both cost \omterm{def:g10:energy-conservation:energy}{energy}. It is the ash of every nuclear fire.
\end{solution}

\begin{solution}{exo:g12:nuclear-energy:14}
Fe: $\Delta m = 26 \times 1.00728 + 30 \times 1.00866 - 55.92066 =
\qty{0.52842}{u}$, $E_b \approx \qty{492}{MeV}$, $E_b/A \approx
\qty{8.79}{MeV}$. U: $\Delta m = 92 \times 1.00728 + 143 \times
1.00866 - 234.99350 = \qty{1.91464}{u}$, $E_b \approx \qty{1784}{MeV}$,
$E_b/A \approx \qty{7.59}{MeV}$. Repacking $235$ \omterm{def:g11:fundamental-interactions:ladder}{nucleons} at
\qty{8.5}{MeV} instead of \qty{7.59}{MeV} gains $235 \times 0.9
\approx \qty{2e2}{MeV}$ --- the $185$--$200$ of the worked example.
\end{solution}

\begin{solution}{exo:g12:nuclear-energy:15}
$\num{5.0e8}/\num{2.82e-12} \approx \num{1.8e20}$ reactions per
second; mass rate $\num{1.8e20} \times 5.029 \times \num{1.66054e-27}
\approx \qty{1.5e-6}{kg/s}$, i.e.\ about \qty{0.13}{kg} per day ---
against \num{1.4e6} \omterm{def:g10:orders-of-magnitude:unit}{kilograms} of coal: ten million to one.
\end{solution}

\begin{solution}{pb:g12:nuclear-energy:1}
\textbf{1.} $P_{\text{th}} = \num{1.0e9}/0.33 \approx \qty{3.0}{GW}$.

\textbf{2.} $200 \times \num{1.60e-13} = \qty{3.2e-11}{J}$.

\textbf{3.} $\num{3.0e9}/\num{3.2e-11} \approx \num{9.4e19}$ \omterm{def:g12:nuclear-energy:fission}{fissions}
per second.

\textbf{4.} $m_{\mathrm U} = 235 \times \num{1.66054e-27} =
\qty{3.90e-25}{kg}$; $\num{9.4e19} \times \num{3.90e-25} \approx
\qty{3.7e-5}{kg/s}$ --- some forty micrograms a second.

\textbf{5.} $\num{3.7e-5} \times \num{3.16e7} \approx \qty{1.2e3}{kg}$:
about $1.2$ tonnes of uranium-235 a year.

\textbf{6.} $\num{1.2e3}/0.040 \approx \qty{2.9e4}{kg}$ --- some $29$
tonnes of fuel: one heavy truck, once a year.

\textbf{7.} $\num{3.0e9} \times \num{3.16e7} \approx \qty{9.5e16}{J}$.

\textbf{8.} $\num{9.5e16}/\num{3.0e7} \approx \qty{3.2e9}{kg}$: $3.2$
million tonnes of coal.

\textbf{9.} $\num{3.2e9}/\num{1.2e3} \approx \num{2.7e6}$ --- the
$\num{8.2e13}/\num{3.0e7}$ of the ladder, as it must be.

\textbf{10.} $\num{3.2e9}/\num{3.0e6} \approx 1060$ trains a year,
roughly three a day: a coal train at dawn, noon and dusk for ever,
against one truck each spring.

\textbf{11.} $17.6 \times \num{1.60e-13} = \qty{2.82e-12}{J}$.

\textbf{12.} One pair weighs $5.029 \times \num{1.66054e-27} =
\qty{8.35e-27}{kg}$: $\num{2.82e-12}/\num{8.35e-27} \approx
\qty{3.4e14}{J/kg}$ --- the ladder's \omterm{def:g12:nuclear-energy:fusion}{fusion} rung.

\textbf{13.} $\num{9.5e16}/\num{3.4e14} \approx \qty{2.8e2}{kg}$ of
D--T; deuterium share $280 \times 2.014/5.029 \approx \qty{1.1e2}{kg}$.

\textbf{14.} $\num{1.1e2}/\num{3.3e-2} \approx \qty{3.4e3}{m^3}$ of
seawater --- about a pool and a half for the whole city's year.

\textbf{15.} Nothing multiplies: no neutron lights the next \omterm{def:g12:nuclear-energy:fusion}{fusion},
and losing confinement cools the \omterm{def:g12:nuclear-energy:barrier}{plasma} and stops it within seconds.
The hard part is the opposite --- holding \qty{1.5e8}{K} together long
enough to burn.

\textbf{16.} $\num{3.9e26}/\num{9.00e16} \approx \qty{4.3e9}{kg/s}$:
four million tonnes a second.

\textbf{17.} $\num{3.0e9}/\num{9.00e16} \approx \qty{3.3e-8}{kg/s}$,
i.e.\ $\num{3.3e-8} \times \num{3.16e7} \approx \qty{1.0}{kg}$ a year
--- and yes: $\qty{1.2e3}{kg} \times 0.00085 \approx 1$, $\num{3.2e9}
\times \num{3.3e-10} \approx 1$, $280 \times 0.00375 \approx 1$. Every
fuel surrenders the same \omterm{def:g10:orders-of-magnitude:unit}{kilogram}; they differ only in how much cargo
must carry it.

\textbf{18.} $\num{4.3e9} \times \num{3.16e7} \times \num{4.5e9}
\approx \qty{6.1e26}{kg}$ --- only $\num{6.1e26}/\num{2.0e30} \approx
\num{3e-4}$, three parts in ten thousand.

\textbf{19.} Convertible mass $0.0007 \times \num{2.0e30} =
\qty{1.4e27}{kg}$, worth $\num{1.4e27} \times \num{9.00e16} \approx
\qty{1.3e44}{J}$; at \qty{3.9e26}{W} that lasts $\num{1.3e44}/
\num{3.9e26} \approx \qty{3.2e17}{s} \approx \num{1.0e10}$ years ---
ten billion, of which some five billion remain.

\textbf{20.} One year of the city: $3.2$ million tonnes of coal, or
$\num{1.2e3}$ \omterm{def:g10:orders-of-magnitude:unit}{kilograms} of uranium-235, or $\num{2.8e2}$ \omterm{def:g10:orders-of-magnitude:unit}{kilograms} of
D--T drawn from a pool and a half of seawater. Under every ledger the
same entry: one \omterm{def:g10:orders-of-magnitude:unit}{kilogram} of mass, gone. $E = mc^2$ does not care which
fuel carries it --- only how big a truck it takes.
\end{solution}
